\documentclass[a4paper,12pt]{article}
\usepackage[T2A]{fontenc}
\usepackage[utf8]{inputenc}
\usepackage[russian]{babel}
\usepackage{amsmath,amssymb,mathtools}
\usepackage{PTSans}
\renewcommand{\familydefault}{\sfdefault}
\usepackage{icomma}
\usepackage[a4paper,margin=20mm,headsep=7mm,footskip=12mm]{geometry}
\usepackage{microtype}
\usepackage{tikz}
\usetikzlibrary{calc,arrows.meta,positioning,shapes.geometric}
\usepackage{tabularx,booktabs,enumitem}
\usepackage{hyperref}
\usepackage{parskip}
\usepackage{fancyhdr}
\usepackage{setspace}
\usepackage{xcolor}

\definecolor{accent}{HTML}{2F6F73}
\definecolor{darkaccent}{HTML}{19484B}
\definecolor{textgray}{HTML}{2E3336}
\definecolor{softgray}{HTML}{687579}
\definecolor{rulegray}{HTML}{CAD5D5}
\hypersetup{colorlinks=true,linkcolor=darkaccent,urlcolor=darkaccent}
\color{textgray}
\onehalfspacing
\setlength{\parindent}{0pt}
\setlist[itemize]{leftmargin=6mm,itemsep=1mm,topsep=1.2mm}
\setlist[enumerate]{leftmargin=7mm,itemsep=2mm,topsep=2mm}
\setlength{\headheight}{15pt}
\setlength{\tabcolsep}{5pt}
\renewcommand{\arraystretch}{1.21}
\pagestyle{fancy}
\fancyhf{}
\fancyhead[L]{\small\color{softgray}Текстовые задачи: работа и растворы}
\fancyhead[R]{\small\color{softgray}08.10.26}
\fancyfoot[C]{\small\color{softgray}\thepage}
\renewcommand{\headrulewidth}{0pt}
\newcommand{\sectiontitle}[1]{%
  \vspace{2mm}{\Large\bfseries\color{darkaccent}#1}\par
  \vspace{1.2mm}{\color{rulegray}\hrule}\vspace{3mm}}
\newcommand{\key}[1]{\textbf{\color{darkaccent}#1}}
\newcommand{\formula}[1]{\[\displaystyle #1\]}
\newcommand{\unit}[1]{\,\text{#1}}

\begin{document}
\thispagestyle{empty}
\begin{center}
\vspace*{31mm}
{\small\color{softgray}ЧЕК-ЛИСТ}\par\vspace{7mm}
{\Huge\bfseries\color{darkaccent}Текстовые задачи\\[2mm]на работу и растворы}\par
\vspace{7mm}
{\Large Производительность, системы, смеси}\par
\vspace{18mm}
{\large 08.10.26}\par
\vfill
{\large\bfseries Лёвин Артём Александрович}\par\vspace{2mm}
{\large эксклюзивно для Матвея Горбачёва}\par
\vspace{20mm}
\end{center}
\begin{tikzpicture}[remember picture,overlay]
\draw[accent,line width=1.2pt]
  ($(current page.south west)+(15mm,18mm)$)
  .. controls ($(current page.south west)+(68mm,32mm)$) and
  ($(current page.south east)+(-71mm,7mm)$) ..
  ($(current page.south east)+(-15mm,20mm)$);
\draw[rulegray,line width=0.7pt]
  ($(current page.south west)+(15mm,13mm)$)
  .. controls ($(current page.south west)+(71mm,25mm)$) and
  ($(current page.south east)+(-67mm,3mm)$) ..
  ($(current page.south east)+(-15mm,15mm)$);
\end{tikzpicture}
\clearpage

\sectiontitle{1. Единый табличный метод}

У двух тем общая логика: выбрать величину, которая \emph{сохраняется или складывается}, заполнить таблицу и составить уравнение по условию.

\begin{center}
\begin{tabularx}{\linewidth}{@{}p{36mm}XXX@{}}
\toprule
Тип задачи & Первая величина & Вторая величина & Результат \\
\midrule
Работа & Производительность $p$ & Время $t$ & Объём работы $A=pt$ \\
Смеси & Масса смеси $m$ & Доля вещества $c$ & Масса вещества $q=mc$ \\
\bottomrule
\end{tabularx}
\end{center}

\key{Алгоритм.} (1) Определите вопрос задачи и выберите неизвестное; (2) постройте таблицу; (3) заполните недостающий столбец по формуле; (4) выразите словесное условие уравнением; (5) решите и проверьте единицы, ограничения и смысл ответа.

\sectiontitle{2. Совместная работа: $A=pt$}

Производительность показывает, какая доля заказа выполняется за единицу времени. Если целый заказ принять за $1$, то исполнитель, выполняющий его за $T$ часов, имеет производительность $1/T$ заказа в час.
\formula{A=pt,\qquad t=\frac{A}{p},\qquad p=\frac{A}{t}.}

Если исполнители одновременно выполняют \emph{одну} работу, их производительности складываются:
\formula{p_{\text{общ}}=p_1+p_2,\qquad \frac1T=\frac1{T_1}+\frac1{T_2}.}

\key{Пример.} Первый мастер выполняет заказ за $3$ часа, второй --- за $6$ часов. Вместе за час они выполнят $\frac13+\frac16=\frac12$ заказа. Поэтому $T=2\unit{ч}$.

\sectiontitle{3. Время как отношение результата к скорости работы}

Петя отвечает на $6$ вопросов за час, Ваня --- на $7$. Одинаковый тест Петя закончил на $20$ минут позже; начали они одновременно. Пусть в тесте $x$ вопросов.
\begin{center}
\begin{tabularx}{0.88\linewidth}{@{}lXXX@{}}
\toprule
& Производительность & Время & Результат\\
\midrule
Петя & $6$ & $x/6$ & $x$\\
Ваня & $7$ & $x/7$ & $x$\\
\bottomrule
\end{tabularx}
\end{center}

Разность времён выражается \emph{в часах}: $20\unit{мин}=\frac13\unit{ч}$.
\formula{\frac{x}{6}-\frac{x}{7}=\frac13 \quad\Longrightarrow\quad x=14.}

\clearpage

\sectiontitle{4. Три исполнителя: что именно нужно найти?}

Игорь и Паша выполняют работу за $3$ часа, Паша и Володя --- за $6$ часов, Володя и Игорь --- за $4$ часа. Обозначим личные производительности через $a,b,c$.
\formula{\begin{cases}
a+b=\frac13,\\ b+c=\frac16,\\ c+a=\frac14.
\end{cases}}

\key{Если нужна общая производительность}, достаточно сложить три уравнения: каждый исполнитель встретится два раза.
\formula{2(a+b+c)=\frac13+\frac16+\frac14=\frac34,\qquad
 a+b+c=\frac38.}
Совместное время втроём: $T=1/(3/8)=\frac83\unit{ч}=2\unit{ч}\ 40\unit{мин}$.

\key{Если нужны личные производительности}, найдите каждую неизвестную, например:
\formula{a=\frac{(a+b)+(a+c)-(b+c)}{2}.}
Затем переведите производительности во времена по формуле $T_i=1/p_i$. Не путайте производительность и время выполнения заказа.

\sectiontitle{5. Трубы и изменение состава бригад}

\key{Трубы.} Первая пропускает $x$ л/мин, вторая --- на $12$ л/мин больше. Для бака объёмом $160$ л первая тратит на $12$ мин больше:
\formula{\frac{160}{x}-\frac{160}{x+12}=12,\qquad x>0.}
Решение даёт $x=8\unit{л/мин}$; отрицательный корень исключается по смыслу.

\key{Бригады.} Два одинаковых заказа начали одновременно $16$ и $25$ работников одной квалификации. Через $7$ дней из второй бригады в первую перешли $8$ работников. Пусть $p>0$ --- производительность одного работника, а $t$ --- число дней \emph{после перехода}. При одновременном завершении объёмы работ равны:
\formula{7\cdot16p+24pt=7\cdot25p+17pt.}
Сокращаем на $p$, получаем $t=9$. \key{Ответ на вопрос} о полном времени: $7+9=16\unit{дней}$.

\key{Типичная ошибка.} В задачах с изменением условий обязательно разделяйте время на этапы «до» и «после» и возвращайтесь к полной длительности.
\clearpage

\sectiontitle{6. Смеси и растворы: сохраняется масса вещества}

В задачах на растворы массовую долю удобно обозначать $c$ и записывать \emph{десятичной дробью}: $15\%=0,15$.
\formula{c=\frac{q}{m},\qquad q=mc,\qquad
 c_{\text{итог}}=\frac{m_1c_1+m_2c_2}{m_1+m_2}.}
Здесь $m$ --- масса всего раствора, $q$ --- масса растворённого вещества. При смешивании без потерь складываются \emph{массы растворов и массы вещества}; доли напрямую не складываются.

\key{Пример.} Смешали $4$ части $15\%$ раствора и $6$ частей $25\%$ раствора того же вещества.
\begin{center}
\begin{tabularx}{0.95\linewidth}{@{}lXXX@{}}
\toprule
Раствор & Масса $m$ & Доля $c$ & Вещество $q=mc$\\
\midrule
Первый & $4$ & $0,15$ & $0,6$\\
Второй & $6$ & $0,25$ & $1,5$\\
После смешивания & $10$ & $x$ & $10x$\\
\bottomrule
\end{tabularx}
\end{center}
\formula{10x=0,6+1,5\quad\Longrightarrow\quad x=0,21=21\%.}

\key{Важное допущение.} Если в условии заданы \emph{литры}, для использования таблицы массовых долей необходимо знать плотности. В учебной задаче выше можно считать объёмы пропорциональными массам (равные плотности); без такого допущения данных об объёмах недостаточно.

\key{Равные массы.} Если смешаны одинаковые массы $15\%$ и $19\%$ растворов, то
\formula{2mc=0,15m+0,19m\quad\Longrightarrow\quad c=0,17=17\%.}
Это среднее арифметическое долей \emph{именно при равных массах}.

\key{Чистая вода} имеет массовую долю растворённого вещества $0$; при её добавлении масса вещества остаётся неизменной, а общая масса увеличивается.

\sectiontitle{7. Самопроверка перед тренировкой}
\begin{itemize}[label=$\square$,itemsep=1.5mm]
\item Умею заполнить таблицу «производительность --- время --- результат».
\item Умею составить уравнение по разности времён и перевести минуты в часы.
\item Умею найти сумму производительностей трёх исполнителей по работе парами.
\item Разделяю работу бригад на этапы до изменения состава и после него.
\item Умею составить баланс массы вещества при смешивании растворов.
\item Проверяю единицы, положительность времени и допустимость концентрации.
\end{itemize}

\clearpage
\thispagestyle{plain}
\begin{center}
{\Large\bfseries\color{darkaccent}Алгоритм: от текста задачи к уравнению}
\end{center}
\vspace{5mm}
\begin{center}
\begin{tikzpicture}[
  x=1cm,y=1cm,
  every node/.style={font=\small,align=center},
  start/.style={ellipse,draw=accent,line width=1pt,minimum width=5cm,minimum height=0.95cm,inner sep=2mm},
  action/.style={rounded corners=2mm,rectangle,draw=rulegray,line width=0.95pt,text width=11.0cm,minimum height=1.08cm,inner sep=2.5mm},
  branch/.style={rounded corners=2mm,rectangle,draw=rulegray,line width=0.95pt,text width=5.3cm,minimum height=1.4cm,inner sep=2mm},
  decision/.style={diamond,aspect=2.75,draw=accent,line width=1pt,text width=3.4cm,inner sep=1.5mm},
  arrow/.style={-{Stealth[length=2.4mm]},line width=0.9pt,draw=softgray}
]
\node[start] (s) at (0,0) {Прочитать условие и вопрос};
\node[action] (q) at (0,-1.55) {Выбрать неизвестное; согласовать единицы};
\node[decision] (d) at (0,-3.45) {Какая задача?};
\node[branch] (w) at (-4.05,-5.70) {Работа: таблица $p$, $t$, $A$\\$A=pt$};
\node[branch] (m) at (4.05,-5.70) {Растворы: таблица $m$, $c$, $q$\\$q=mc$};
\node[action] (f) at (0,-8.15) {Заполнить таблицу: известные данные, неизвестные величины, результаты вычислений};
\node[action] (r) at (0,-10.05) {Найти словесную связь: равенство объёмов, разность времени, сумма производительностей или масса вещества};
\node[action] (eq) at (0,-11.95) {Составить уравнение или систему и решить};
\node[action] (v) at (0,-13.85) {Проверить ограничения, единицы и смысл результата; ответить на вопрос};
\node[start] (e) at (0,-15.45) {Ответ};
\draw[arrow] (s)--(q);
\draw[arrow] (q)--(d);
\draw[arrow] (d)--node[above left,near end]{работа}(w);
\draw[arrow] (d)--node[above right,near end]{растворы}(m);
\draw[arrow] (w.south)--++(0,-0.48)-|(f.north west);
\draw[arrow] (m.south)--++(0,-0.48)-|(f.north east);
\draw[arrow] (f)--(r);
\draw[arrow] (r)--(eq);
\draw[arrow] (eq)--(v);
\draw[arrow] (v)--(e);
\end{tikzpicture}
\end{center}
\clearpage

\sectiontitle{8. Тренировочный блок --- Путь А}
\textcolor{softgray}{Для каждой задачи составьте таблицу и уравнение. Запишите решение и ответ. Задания расположены по нарастающей сложности внутри каждого типа.}

\key{Работа и производительность}
\begin{enumerate}[label=\textbf{\arabic*.},itemsep=2.7mm]
\item Один мастер выполняет заказ за $4$ ч, второй --- за $12$ ч. За какое время они выполнят этот заказ вместе?
\item Петя отвечает на $8$ вопросов в час, Ваня --- на $10$ вопросов в час. Они одновременно начали одинаковые тесты; Петя закончил на $18$ мин позже. Сколько вопросов в каждом тесте?
\item Первая и вторая бригады выполняют заказ вместе за $4$ ч, вторая и третья --- за $4$ ч, первая и третья --- за $8$ ч. За сколько часов три бригады выполнят заказ вместе?
\item Первый и второй мастера выполняют заказ за $4$ ч, второй и третий --- за $6$ ч, третий и первый --- за $4$ ч. За сколько часов \emph{каждый} выполнит заказ отдельно?
\item Первая труба пропускает на $5$ л воды в минуту меньше второй. Бак объёмом $150$ л первая заполняет на $5$ мин дольше второй. Найдите производительность первой трубы.
\item Две бригады из $14$ и $20$ работников одинаковой квалификации одновременно начали выполнять одинаковые заказы. Через $6$ дней из второй бригады в первую перешли $5$ работников. Оба заказа завершили одновременно. Сколько дней длилась работа?
\end{enumerate}

\vspace{2mm}\key{Смеси и растворы}
\begin{enumerate}[label=\textbf{\arabic*.},start=7,itemsep=2.7mm]
\item Смешали $3$ кг $10\%$ раствора и $7$ кг $30\%$ раствора того же вещества. Найдите массовую долю вещества в смеси.
\item Смешали равные массы $16\%$ и $28\%$ растворов одного вещества. Найдите концентрацию итогового раствора.
\item К $8$ кг $30\%$ раствора добавили некоторое количество чистой воды. Получился $20\%$ раствор. Сколько килограммов воды добавили?
\item Из $10\%$ и $25\%$ растворов одного вещества нужно получить $18$ кг $20\%$ раствора. Сколько килограммов каждого раствора потребуется?
\end{enumerate}

\clearpage
\sectiontitle{9. Ответы для самопроверки}
\begin{center}
\renewcommand{\arraystretch}{1.5}
\begin{tabularx}{\linewidth}{@{}p{16mm}X@{}}
\toprule
Номер & Ответ и контрольное соотношение\\
\midrule
1 & $3$ ч; $\frac1T=\frac14+\frac1{12}=\frac13$.\\
2 & $12$ вопросов; $\frac{x}{8}-\frac{x}{10}=\frac{18}{60}$.\\
3 & $3$ ч $12$ мин; $\frac1T=\frac12\left(\frac14+\frac14+\frac18\right)=\frac5{16}$.\\
4 & Первый --- $6$ ч, второй --- $12$ ч, третий --- $12$ ч.\\
5 & $10$ л/мин; $\frac{150}{x}-\frac{150}{x+5}=5$.\\
6 & $15$ дней; $6\cdot14+19t=6\cdot20+15t$, $t=9$.\\
7 & $24\%$; $\frac{3\cdot0,10+7\cdot0,30}{10}=0,24$.\\
8 & $22\%$; $\frac{0,16+0,28}{2}=0,22$.\\
9 & $4$ кг; $8\cdot0,30=(8+x)\cdot0,20$.\\
10 & $6$ кг $10\%$ раствора и $12$ кг $25\%$ раствора.\\
\bottomrule
\end{tabularx}
\end{center}

\vspace{5mm}
\key{Как себя проверить.} Для задач о работе подставьте найденное время обратно в выражение объёма работы. Для растворов убедитесь, что суммарная масса вещества до и после смешивания одинакова. В каждом случае ответ должен иметь смысл: время и производительности положительны, массовая доля находится между $0\%$ и $100\%$.

\end{document}
